\section{Исследование и построение решения задачи}
\label{sec:Chapter3} \index{Chapter3}

%Требуется разбить большую задачу, описанную в постановке, на более мелкие
%подзадачи. Процесс декомпозиции следует продолжать до тех пор, пока подзадачи
%не станут достаточно простыми для решения непосредственно. Это может быть
%достигнуто, например, путем проведения эксперимента, доказательства теоремы
%или поиска готового решения.

\subsection{Введение в Linux memory management}

Перед тем, как перейти к описанию разработанного механизма, необходимо рассмотреть основы управления пользовательской памятью в ядре Linux. В следующих разделах кратко рассматриваются виртуальный и физический уровни памяти, а также NUMA-балансировка, необходимая для дальнейшего изложения.

\subsubsection{Физический уровень}

В современных вычислительных системах управление памятью осуществляется с использованием механизма виртуальной памяти, при котором каждый процесс работает с собственным виртуальным адресным пространством. Преобразование виртуальных адресов в физические выполняется аппаратным блоком управления памятью (MMU) с использованием иерархии таблиц трансляций. В частности, в архитектуре arm64 используются две иерархии таблиц трансляций: одна для \textit{user-space}, вторая~--- для \textit{kernel-space} (см рис.~\ref{fig:TTBR_arm64}). Физические адреса верхнеуровневых PGD-таблиц хранятся в регистрах \texttt{TTBR0} и \texttt{TTBR1} соответственно. В дальнейшем нас будет интересовать \textit{user-space} иерархия.

Ключевым понятием в механизме трансляции адресов является \textit{базовый размер страницы (page size)}~--- он определяет минимальную единицу отображения виртуальной памяти на физическую. Таким образом, \textit{виртуальные страницы} базового размера отображаются на \textit{физические страницы (page frames)}, того же размера, а сами страницы выровнены по его границе. Конкретное значение размера задается архитектурой процессора, ядро Linux выбирает один из поддерживаемых вариантов и использует его во всей системе. Выбранный размер влияет на структуру таблиц трансляций, включая число уровней и разбиение виртуального адреса на поля индексации. arm64 поддерживает размеры страниц 4K, 16K, 64K, на рисунке~\ref{fig:TTBR_arm64} изображена структура для 4K страниц, т.к. она наиболее распространенная.

В рамках данной работы помимо 4K страниц интерес представляют ещё и 64K страницы, потому что они часто используются в серверных системах~--- целевых для работы. Но в Linux средства взаимодействия с таблицами трансляций реализованы с помощью специальных трюков, основанных на условной компиляции, таким образом, что код, работающий для одного базового размера страницы, работает должным образом и для другого (один из таких трюков~--- \textit{page table folding}\cite{page_table_folding}). Поэтому для разработки и реализации механизма значение размера страницы не играет существенной роли, и в дальнейшем мы сосредоточимся на 4K страницах; для 64K страниц все будет работать аналогично.

Важно отметить, что в arm64 каждая таблица трансляций (на любом уровне иерархии) занимает ровно одну страницу~--- это упрощает управление памятью. Также записи таблиц трансляций содержат атрибуты доступа (\textit{read}/\textit{write}/\textit{execute}) отображаемой памяти, при этом страница является наименьшей неделимой единицей, для которой эти атрибуты могут быть заданы. Таким образом, вся физическая память разбивается логически на физические страницы, и Linux для каждой такой страницы хранит \texttt{page}-дескриптор~--- экземпляр структуры \texttt{page}, в котором хранится вся метаинформация, связанная со страницей (см рис.~\ref{fig:struct_page}).

\begin{figure}[H]
    \centering
    \includegraphics[width=1\textwidth]{TTBR_arm64.pdf}
    \caption{Структура таблиц трансляций arm64 для 4K страниц}
    \label{fig:TTBR_arm64}
\end{figure}

\subsubsection{Виртуальный уровень}

Виртуальное адресное пространство любого процесса в Linux состоит из непересекающихся \textit{областей виртуальной памяти (virtual memory areas)} (см рис.~\ref{fig:vma}). Каждая такая область является полуинтервалом (границы которого выровнены по размеру страницы) виртуальных адресов процесса, обладающим единым набором атрибутов и характеристик отображения, таких как права доступа (\textit{read}/\textit{write}/\textit{execute}/...), тип отображения (\textit{file-backed}/\textit{анонимный}), NUMA-политика размещения и прочие. Область описывает логически однородный участок адресного пространства и используется ядром для управления отображением виртуальной памяти на физическую. В ядре Linux каждая область представляется структурой \texttt{vm\_area\_struct}, содержащей всю необходимую информацию о ней.

\begin{figure}[H]
    \centering
    \includegraphics[width=1\textwidth]{struct_page.pdf}
    \caption{struct page}
    \label{fig:struct_page}
\end{figure}

Выделение памяти в виртуальном адресном пространстве процесса происходит в два этапа (см рис.~\ref{fig:mem_alloc}). На первом этапе формируется виртуальная область памяти: ядро резервирует диапазон виртуальных адресов и создаёт соответствующую структуру \texttt{vm\_area\_struct}. Для этого в Linux предусмотрен системный вызов \texttt{mmap}, позволяющий пользовательским программам запрашивать создание таких отображений.

\begin{figure}[H]
    \centering
    \includegraphics[width=1\textwidth]{vma.pdf}
    \caption{Виртуальное адресное пространство}
    \label{fig:vma}
\end{figure}

На втором этапе, при обращении процесса к ранее зарезервированным адресам, возникают \textit{ошибки страниц (page faults)}, в результате которых ядро выделяет необходимые физические страницы и устанавливает соответствующие отображения в таблицах трансляций. При этом, помимо страниц, непосредственно содержащих данные, при необходимости также создаются недостающие структуры таблиц трансляций. Информация, содержащаяся в \texttt{vm\_area\_struct} (например, указатели на функции обратного вызова), позволяет ядру определить, каким именно образом следует разрешить конкретную ошибку страницы.

Второй этап также может быть выполнен немедленно после первого. Для этого при вызове \texttt{mmap}\cite{mmap} может быть указан флаг \texttt{MAP\_POPULATE}, в результате чего ядро заранее выделит необходимые физические страницы и сформирует соответствующие отображения для всего диапазона адресов. В этом случае последующие обращения по данным адресам не приведут к возникновению ошибок страниц. Такой подход в ряде случаев может быть полезен.

\begin{figure}[H]
    \centering
    \includegraphics[width=1\textwidth]{mem_alloc.pdf}
    \caption{Процесс выделения памяти в пространстве пользователя}
    \label{fig:mem_alloc}
\end{figure}

\subsubsection{NUMA-балансировщик}

Механизм автоматической NUMA-балансировки уже был рассмотрен ранее (см главу~\ref{sec:Chapter2_NUMA_balancer}). В рамках данной работы нас будет интересовать его более низкоуровневая организация, в частности, взаимодействие с внутренними структурами ядра, такими как \texttt{vm\_area\_struct} и \texttt{page}.

В упрощённом виде работа балансировщика на каждой итерации сканирования может быть представлена как обход виртуального адресного пространства процесса. Сначала ядро последовательно проходит по структурам \texttt{vm\_area\_struct}, описывающим области памяти процесса, после чего для выбранных областей выполняется обход соответствующих им записей PTE-таблиц трансляций. В ходе этого обхода часть страниц помечается специальным образом (через \textit{invalid} бит в соответствующей PTE-записи), что приводит к генерации так называемых \textit{NUMA hinting faults} при последующих обращениях.

Возникающие таким образом ошибки страниц используются ядром для сбора информации о фактической локальности доступа к памяти, которая затем применяется при принятии решений о миграции страниц между NUMA-нодами.

Общая схема работы NUMA-балансировщика с учётом описанных деталей представлена на рис.~\ref{fig:numa_balancer}.

\begin{figure}[ht]
    \centering
    \includegraphics[width=1\textwidth]{numa_balancer.pdf}
    \caption{Схема работы NUMA-балансировщика}
    \label{fig:numa_balancer}
\end{figure}

\subsection{Дизайн механизма}

Перейдём к описанию предлагаемого механизма репликации. Как было отмечено во введении, его основная идея заключается в создании локальных копий таблиц трансляций и данных для различных NUMA-нод с целью повышения локальности доступа.

В рамках данной работы репликация применяется только к \textit{private read-only}/ \textit{executable} данным. Это обусловлено тем, что репликация \textit{writable} данных требует поддержки когерентности между копиями и может приводить к существенным накладным расходам, нивелируя потенциальный выигрыш от повышения локальности. Аналогичное ограничение распространяется и на \textit{shared read-only}/\textit{executable} данные: несмотря на отсутствие возможности записи со стороны рассматриваемого процесса, такие отображения могут быть изменены другими процессами или ядром операционной системы, а эти изменения должны оставаться наблюдаемыми. В результате их репликация также требует механизмов поддержания согласованности между копиями.

Таблицы трансляций, стоит сказать, как и \textit{writable} данные, тоже являются изменяемыми структурами. Однако их свойства принципиально отличаются:
\begin{enumerate}
    \item Во-первых, изменения таблиц страниц в типичном случае происходят относительно редко по сравнению с интенсивностью обращений к данным, ограничиваясь в основном этапами инициализации и завершения, а также отдельными событиями управления памятью.
    \item Во-вторых, все модификации таблиц выполняются исключительно в пространстве ядра, что позволяет отслеживать их и интегрировать соответствующую логику в реализуемый механизм. В то время как изменения \textit{writable} данных в общем случае происходят без участия ядра, напрямую аппаратурой, что сильно усложняет обеспечение согласованности их копий.
\end{enumerate}
И, наконец, репликация таблиц трансляций является необходимым условием репликации самих данных, поскольку именно они задают отображение виртуальных страниц на физические и обеспечивают доступ к локальным копиям на различных NUMA-нодах.

На рисунке~\ref{fig:replicated_data} изображены реплицированные таблицы трансляций. Такая организация возможна благодаря тому, что каждый CPU использует собственный регистр \texttt{TTBR0}, в который может быть записан адрес PGD-таблицы, размещённой на локальной NUMA-ноде. Если же реплицированы будут все уровни таблиц трансляций, появится возможность реплицировать и сами данные: каждая копия данных будет связана с соответствующей записью в локальной копии PTE-таблицы (см RO (\textit{read-only}) данные на рис.~\ref{fig:replicated_data}). Таким образом, одной виртуальной странице будут соответствовать несколько физических, каждая из которых будет доступна потокам, выполняющимся на CPU в пределах своей NUMA-ноды.

Однако данные можно и не реплицировать~--- в таком случае все копии PTE-записей (соответствующие одной и той же виртуальной странице) будут указывать на одну и ту же физическую страницу, расположенную на одной из NUMA-нод. При этом доступ к данной странице сохранится с любой NUMA-ноды и будет осуществляться через соответствующую локальную иерархию таблиц трансляций (см RW (\textit{read-write}) данные на рис.~\ref{fig:replicated_data}). Такой подход является необходимой мерой для \textit{writable} данных, не подлежащих репликации в рамках выбранного дизайна, но может применяться и для других категорий данных. Репликация таблиц трансляций может быть полезна и без репликации данных, выступая средством снижения NUMA-эффекта, обусловленного обращениями к таблицам трансляций, которые также могут размещаться на удалённых NUMA-нодах.

\begin{figure}[H]
    \centering
    \includegraphics[width=1\textwidth]{replicated_data.pdf}
    \caption{Реплицированные таблицы трансляций и данные в NUMA-системе}
    \label{fig:replicated_data}
\end{figure}

Репликация может завершаться не только на уровне данных, но и на одном из промежуточных уровней таблиц трансляций (см рис.~\ref{fig:replicated_tables}). Это связано с тем, что при репликации таблиц для заданного диапазона виртуальных адресов необходимо реплицировать все структуры, лежащие на пути обхода соответствующих записей. При этом отдельные таблицы могут содержать ссылки на дочерние структуры, не относящиеся к рассматриваемому диапазону, и, следовательно, не подлежащие репликации. В таких случаях репликация иерархии обрывается на соответствующем уровне.

В рамках предлагаемого дизайна корневой уровень иерархии, состоящий из единственной PGD-таблицы, всегда реплицирован~--- это упростит дальнейшую работу механизма с таблицами трансляций.

Таким образом, в общем случае репликация распространяется от корня иерархии до некоторого уровня, на котором она может быть прервана. На этом уровне записи таблиц трансляций начинают ссылаться на единые (нереплицированные) структуры следующего уровня, и дальнейшая трансляция адресов выполняется по стандартной схеме с использованием единственного экземпляра соответствующих таблиц и физических страниц.

\begin{figure}[H]
    \centering
    \includegraphics[width=1\textwidth]{replicated_tables.pdf}
    \caption{Частично реплицированные таблицы трансляций}
    \label{fig:replicated_tables}
\end{figure}

\subsubsection{Политики работы и пользовательский интерфейс}
\label{sec:Chapter3_user_API}

Существенной особенностью пользовательской памяти является её динамическая природа: время жизни и характер использования данных определяются пользовательским приложением и могут изменяться во время выполнения программы. При этом ядро не располагает полной информацией о семантике и ожидаемом времени жизни отдельных областей памяти, что затрудняет принятие оптимальных решений о репликации.

Поэтому в предлагаемом дизайне предусмотрено несколько политик работы механизма, определяющих стратегию репликации. Есть отдельные политики репликации таблиц (см таб.~\ref{tab:table_policy}) и отдельные политики репликации данных (см таб.~\ref{tab:data_policy}), итоговое поведение механизма зависит от комбинации двух политик. Все комбинации возможны, кроме (\texttt{None}, \texttt{Lazy}) и (\texttt{None}, \texttt{All}), потому что репликация данных невозможна без репликации таблиц.

\begin{table}[H]
    \centering
    \begin{tabular}{
        |>{\centering\arraybackslash}m{2cm}
        |>{\centering\arraybackslash}m{13cm}|}
        \hline
        \textbf{Политика} & \textbf{Описание} \\ \hline
        None    & Таблицы не реплицируются \\ \hline
        Minimal & Таблицы для \textit{private read-only} данных реплицируются \\ \hline
        All     & Все таблицы реплицируются \\ \hline
    \end{tabular}
    \caption{Политики репликации таблиц трансляций}
    \label{tab:table_policy}
\end{table}
\begin{table}[H]
    \centering
        \begin{tabular}{
        |>{\centering\arraybackslash}m{2cm}
        |>{\centering\arraybackslash}m{13cm}|}
        \hline
        \textbf{Политика} & \textbf{Описание} \\ \hline
        None & Данные не реплицируются \\ \hline
        Lazy & Горячие* данные реплицируются NUMA-балансировщиком в ленивой манере \\ \hline
        All  & Данные реплицируются немедленно \\ \hline
        \multicolumn{2}{l}{\footnotesize * Горячие данные~--- данные, помеченные NUMA-балансировщиком для \textit{NUMA hinting faults}} \\
    \end{tabular}
    \caption{Политики репликации данных}
    \label{tab:data_policy}
\end{table}

В случае таблиц трансляций политики определяют набор виртуальных областей памяти, для которых допускается репликация таблиц. В зависимости от характера приложения более эффективной может оказаться как репликация всех таблиц трансляций, так и только таблиц для \textit{private read-only} областей.

Для репликации данных политики имеют иной смысл. Набор потенциально реплицируемых областей памяти фиксирован и ограничивается \textit{private read-only} отображениями, тогда как сами политики определяют стратегию выполнения репликации данных: выполнять её немедленно для всех подходящих страниц или только для страниц, активно используемых на разных NUMA-нодах во время выполнения приложения. Политика \texttt{All} ориентирована на максимально агрессивное использование механизма репликации и наиболее эффективна в приложениях с ярко выраженным NUMA-эффектом, где большая часть памяти активно используется на разных NUMA-нодах. Политика \texttt{Lazy}, напротив, стремится реплицировать только действительно горячие данные, избегая лишних накладных расходов на создание и поддержание реплик. Такой подход может быть предпочтительнее для приложений с менее предсказуемым или неоднородным характером доступа к памяти.

Механизм политик также представляет интерес с исследовательской точки зрения, поскольку позволяет сравнить влияние различных вариантов поведения на производительность тестовых программ и тем самым получить более полное представление о работе системы. А также, наличие политик повысит гибкость автоматизации: они делают механизм более модульным и вариативным, что позволит подбирать наиболее эффективную стратегию его применения.

Пользователь может установить политики двумя способами: либо для отдельного процесса, либо для группы процессов. Для этого реализацией предоставляются специальные интерфейсы в \textit{procfs}\cite{procfs} и \textit{memory cgroup fs}\cite{cgroup} соответственно. Интерфейс \textit{cgroup} примечателен тем, что все дочерние процессы наследуют \textit{memory cgroup} родительского процесса, а вместе с ним~--- и установленные политики. Это позволяет, задав нужные политики, например, для командной оболочки или управляющего скрипта, автоматически распространять их на все запускаемые из-под них процессы и создаваемые ими потоки.

Поддержаны \textit{runtime}-переходы между всеми допустимыми комбинациями политик, но только для конкретных процессов, не для групп. Это означает, что смена политик возможна через \textit{procfs}, а также при смене процессом \textit{memory cgroup}. Смена политик на уровне \textit{cgroup}, содержащей ненулевое количество процессов, невозможна.

\subsubsection{Репликация таблиц трансляций}

При рассмотрении вопросов, связанных с выделением физических страниц памяти, центральную роль в ядре Linux играет функция \texttt{handle\_mm\_fault}. Она вызывается в ситуациях, когда требуется обеспечить доступ к физическим данным, в частности при обработке ошибок страниц и при использовании флага \texttt{MAP\_POPULATE}. На вход функция получает контекст обращения, включающий виртуальный адрес и соответствующий экземпляр \texttt{vm\_area\_struct}.

Поскольку доступ к физическим страницам невозможен без соответствующих отображений, неотъемлемой частью работы \texttt{handle\_mm\_fault} является обеспечение наличия необходимых таблиц трансляций. В случае их отсутствия функция последовательно, слой за слоем, выполняет их аллокацию и инициализацию, формируя корректную иерархию таблиц для заданного виртуального адреса.

Именно в этом месте механизм выполняет репликацию таблиц трансляций, необходимость которой определяется установленной для процесса политикой и характеристиками соответствующего экземпляра \texttt{vm\_area\_struct}.

\begin{itemize}
    \item Если репликация для рассматриваемого слоя не требуется, то возможны следующие ситуации:
    \begin{enumerate}
        \item Дочерняя таблица уже существует. В этом случае никаких дополнительных действий не требуется.
        \item Родительская таблица не реплицирована. Это стандартная ситуация, полностью совпадающая с поведением ядра Linux без модификаций.
        \item Родительская таблица реплицирована (см рис.~\ref{fig:table_replication_1}). В этом случае требуется:
        \begin{enumerate}
            \item использовать общий для всех реплик родительского уровня \textit{page table lock}. В Linux для таблиц уровней PMD и PTE ассоциированы собственные блокировки, позволяющие выполнять сериализацию их модификаций более гранулярно;
            \item синхронизировать значения соответствующих записей во всех репликах родительской таблицы.
        \end{enumerate}
        \begin{figure}[H]
            \centering
            \includegraphics[width=1\textwidth]{table_replication_1.pdf}
            \caption{Аллокация нереплицированной PTE-таблицы}
            \label{fig:table_replication_1}
        \end{figure}
    \end{enumerate}

    \item Если репликация требуется, то реализацией гарантируется, что родительская таблица на предыдущем уровне уже реплицирована. В этом случае возможны следующие варианты:
    \begin{enumerate}
        \item Дочерний уровень уже существует и реплицирован. Дополнительные действия не требуются.
        \item Дочерний уровень отсутствует (см рис.~\ref{fig:table_replication_2}). Тогда необходимо:
        \begin{enumerate}
            \item аллоцировать экземпляры таблицы для каждой NUMA-ноды;
            \item захватить общий \textit{page table lock} для всех реплик родительского уровня;
            \item связать каждый дочерний экземпляр с соответствующей репликой родительской таблицы.
        \end{enumerate}
        \begin{figure}[H]
            \centering
            \includegraphics[width=1\textwidth]{table_replication_2.pdf}
            \caption{Аллокация реплицированной PTE-таблицы}
            \label{fig:table_replication_2}
        \end{figure}
        \item Существует единственная (нереплицированная) таблица дочернего уровня (см рис.~\ref{fig:table_replication_3}). В этом случае требуется:
        \begin{enumerate}
            \item выделить дополнительные экземпляры таблицы для остальных NUMA-нод;
            \item захватить общий \textit{page table lock} для родительских реплик, а также \textit{page table lock} существующей дочерней таблицы (что необходимо, поскольку она уже встроена в иерархию и может быть доступна конкурентным потокам);
            \item скопировать содержимое существующей таблицы в новые реплики;
            \item связать каждую реплику с соответствующей родительской таблицей.
        \end{enumerate}
        \begin{figure}[H]
            \centering
            \includegraphics[width=1\textwidth]{table_replication_3.pdf}
            \caption{Репликация существующей PTE-таблицы}
            \label{fig:table_replication_3}
        \end{figure}
    \end{enumerate}
\end{itemize}

Согласно дизайну, наличие реплицированных таблиц трансляций или возможность их появления в будущем для некоторой области виртуальной памяти (в зависимости от установленной политики и характеристик соответствующего \texttt{vm\_area\_struct}) не накладывают ограничений на механизмы ядра, такие как:
\begin{enumerate}
    \item\label{itm:mm_features} Функциональности \textit{KSM}\cite{KSM}, \textit{page migration}\cite{page_migration}, \textit{reverse mapping}\cite{rmap1,rmap2}, \textit{swap}, \textit{userfaultfd}\cite{userfaultfd}, \textit{idle page tracking}\cite{page_idle}. В большинстве своем они взаимодействуют с физическими страницами данных, точкой соприкосновения которых с таблицами трансляций являются записи PTE-таблиц~--- листовых в иерархии. Поэтому для поддержки реплицированных таблиц в них достаточно лишь операции, модифицирующие PTE-записи, заменить на аналоги, модифицирующие соответствующие записи во всех репликах. Это даже не требует понимания внутренней логики функциональностей.

    \item Системные вызовы \texttt{mmap}\cite{mmap}, \texttt{mlock}\cite{mlock}, \texttt{mprotect}\cite{mprotect}, \texttt{madvise}\cite{madvise}. Они оперируют на уровне виртуальных областей памяти и PTE-записей, поддержка реплицированных таблиц осуществляется по аналогии с функциональностями из пункта~\ref{itm:mm_features}.

    \item Системные вызовы \texttt{munmap}\cite{munmap}, \texttt{mremap}\cite{mremap}, \texttt{clone}\cite{clone}. Помимо прочего, выполняют манипуляции с самими таблицами страниц, включая их создание, удаление и перемещение. Это требует более детальной и аккуратной обработки с учетом репликации, однако соответствующие изменения носят локальный характер и не затрагивают ключевые аспекты дизайна механизма, поэтому их подробное рассмотрение сейчас не требуется.
\end{enumerate}

\subsubsection{Репликация данных}
\label{sec:Chapter3_data_repl}

Репликация данных реализуется отдельной функцией \texttt{phys\_duplicate}, принимающей на вход \texttt{vm\_area\_struct} и диапазон виртуальных адресов внутри неё, данные которого необходимо реплицировать. В отличие от репликации таблиц, она выполняется не автоматически в процессе обработки ошибок страниц, а инициируется явно для заданного диапазона.

Сначала \texttt{phys\_duplicate} обходит иерархию таблиц трансляций в пределах заданного диапазона, обеспечивая при необходимости их репликацию~--- для этого используется \texttt{handle\_mm\_fault}. Таким образом, к моменту достижения листового уровня гарантируется, что структура таблиц корректно подготовлена для размещения реплик данных.

Затем выполняется репликация данных на уровне PTE. Для каждой исходной страницы данных предварительно аллоцируются её копии на всех задействованных NUMA-нодах. Содержимое исходной страницы копируется в каждую из реплик, после чего соответствующие PTE-записи обновляются таким образом, что каждая нода использует свою локальную копию страницы. Исходная страница при этом отвязывается от таблиц трансляций и освобождается.

Важно подчеркнуть, что исходная страница намеренно не переиспользуется в качестве одной из реплик. Это решение обусловлено стремлением сделать реплицированные данные максимально изолированными и унифицированными с точки зрения их дальнейшей обработки:
\begin{enumerate}
    \item Во-первых, исходная страница может принадлежать \textit{page cache} (в случае \textit{file-backed} отображений), что потребовало бы дополнительной логики для корректной интеграции с механизмами кэширования файлов. В предлагаемом подходе вместо этого создаётся новая \textit{анонимная} страница, в которую копируется содержимое исходной. Это позволяет работать с единообразным представлением данных независимо от их исходной природы (\textit{анонимной} или \textit{file-backed}), существенно упрощая реализацию.

    \item Во-вторых, реплицированные страницы данных сознательно выводятся из взаимодействия с подсистемами управления памятью \textit{reverse mapping} и \textit{swap}. В частности, для них не используются такие поля структуры \texttt{page}, как \texttt{mapping}, \texttt{\_mapcount}, \texttt{lru}. Переиспользование исходной страницы, уже вовлечённой в эти подсистемы, потребовало бы сложной и потенциально ошибкоопасной очистки и переинициализации её состояния.

    \item В-третьих, даже \textit{анонимные private} страницы могут быть \textit{разделяемыми} между несколькими процессами, например, в результате выполнения системного вызова \texttt{clone}. В этом случае переиспользование исходной страницы привело бы к ситуации, когда одна и та же физическая страница одновременно рассматривается как реплицированная в контексте одного процесса и как обычная в контексте другого. Это нарушило бы изоляцию механизма репликации и привело бы к некорректному поведению (см главу~\ref{sec:Chapter4_phys_layer}).
\end{enumerate}

Таким образом, создание полностью новых, \textit{анонимных} и изолированных копий страниц позволяет избежать дополнительных сложностей и обеспечивает единообразную модель работы с реплицированными данными как для \textit{анонимной}, так и для \textit{file-backed} памяти.

Взаимодействия механизма репликации данных с другими подсистемами управления памятью не столь прозрачны, как у репликации таблиц трансляций. Ниже перечислены инварианты и дизайнерские решения, связанные с этими взаимодействиями:
\begin{enumerate}
    \item\label{itm:swap_rmap} Реплицированные страницы данных в целях простоты не взаимодействуют с функциональностями \textit{reverse mapping}\cite{rmap1, rmap2} и \textit{swap}~--- для них не используются соответствующие поля \texttt{lru}, \texttt{mapping}, \texttt{\_mapcount} структуры \texttt{page}.

    \item Реплицированные данные не могут оказаться в обработчике ошибок страниц, поскольку:
    \begin{itemize}
        \item они представляют собой \textit{анонимные private read-only} данные, поэтому возникновение \textit{copy-on-write} и \textit{spurious write} ошибок исключено;
        \item они не могут находиться в \textit{non-present} состоянии, т.к. не участвуют в \textit{reverse mapping, swap} (см пункт~\ref{itm:swap_rmap}) и \textit{page migration} (см пункт~\ref{itm:migration}) подсистемах;
        \item в коде механизма NUMA-балансировки добавлена проверка, исключающая обработку реплицированных страниц данных, поэтому \textit{NUMA hinting faults} для них также невозможны. Подробнее взаимодействие NUMA-балансировки и механизма репликации данных рассматривается в главе~\ref{sec:Chapter3_data_repl_policies}.
    \end{itemize}

    \item\label{itm:migration} Функциональности \textit{KSM}\cite{KSM} и \textit{page migration}\cite{page_migration} не поддерживают реплицированные данные. Эти подсистемы сначала выполняют обход структур \texttt{vm\_area\_struct} адресного пространства, а уже затем работают с физическими страницами подходящих областей памяти. В соответствующих местах обхода были добавлены проверки, исключающие обработку реплицируемых виртуальных областей. Теоретически возможна более точечная интеграция на уровне отдельных страниц, однако это существенно усложнило бы реализацию и потребовало бы более глубокого вмешательства в логику указанных подсистем, что выходило за рамки поставленной исследовательской задачи.

    \item В системном вызове \texttt{madvise}\cite{madvise} затрагиваются операции \texttt{MADV\_COLD}, \texttt{MADV\_PAGEOUT}, \texttt{MADV\_FREE} и \texttt{MADV\_DONTNEED}. В первых трёх из них реплицированные страницы данных просто пропускаются. Операции \texttt{MADV\_COLD} и \texttt{MADV\_PAGEOUT} связаны с \textit{reclaim (swap)}, который не поддерживается для реплицированных данных, к тому же их семантика носит исключительно рекомендательный характер и допускает игнорирование со стороны реализации. Аналогично \texttt{MADV\_FREE} также не предоставляет пользователю строгих гарантий освобождения памяти.

    Операция \texttt{MADV\_DONTNEED}, напротив, предполагает немедленное освобождение памяти и потому поддерживается корректно. Она использует тот же механизм освобождения, что и системный вызов \texttt{munmap}, поддержка которого является обязательной, поскольку он играет ключевую роль в жизненном цикле управления памятью пользовательских процессов. Благодаря этому удалось избежать дополнительных усилий на поддержку операции.

    \item Системные вызовы \texttt{mmap}\cite{mmap}, \texttt{mlock}\cite{mlock} и \texttt{mprotect}\cite{mprotect} не требуют специальной поддержки реплицированных данных, поскольку не взаимодействуют напрямую с их внутренним представлением.

    \item Для системных вызовов \texttt{munmap}\cite{munmap} и \texttt{clone}\cite{clone} потребовалась дополнительная обработка, связанная с корректным управлением реплицированными данными.

    \item Системный вызов \texttt{mremap}\cite{mremap} не поддерживается для реплицируемых областей памяти. При попытке работы с такими отображениями он завершается с ошибкой \texttt{-EINVAL}. Такое решение обусловлено сложностями в реализации и крайне редким практическим применением данного вызова.

    \item Функциональности \textit{userfaultfd}\cite{userfaultfd} и \textit{idle page tracking}\cite{page_idle} также не поддерживаются, поскольку в рамках данной работы они не представляли практического интереса.
\end{enumerate}

\subsubsection{Подробнее о политиках репликации данных}
\label{sec:Chapter3_data_repl_policies}

Реализацией поддерживаются две политики репликации данных: \texttt{All} и \texttt{Lazy}. Обе из них задействуют функцию \texttt{phys\_duplicate}, однако отличаются моментом и условиями её вызова.

Политика \texttt{All} является наиболее агрессивной: данные реплицируются сразу при появлении возможности для этого. В реализации это достигается достаточно прямолинейно~--- во всех местах ядра, где появляются или модифицируются страницы, потенциально подлежащие репликации, выполняется вызов \texttt{phys\_duplicate} при условии, что для процесса установлена политика \texttt{All}. Таких точек в реализации всего три:
\begin{enumerate}
    \item Аллокация новой физической страницы в обработчике ошибок страниц (\texttt{handle\_mm\_fault}). В этом случае репликация выполняется сразу после создания страницы. Аналогичное поведение возникает и при принудительном выделении памяти, например, с использованием флага \texttt{MAP\_POPULATE}.

    \item Системный вызов \texttt{mprotect}. Если в результате изменения прав доступа область памяти становится \textit{read-only}, то для уже выделенных под неё физических страниц инициируется репликация.

    \item Смена политики репликации данных процесса. При переключении на политику \texttt{All} выполняется обход адресного пространства процесса с последующей репликацией всех подходящих страниц данных.
\end{enumerate}

Политика \texttt{Lazy}, напротив, ориентирована на более избирательную репликацию данных и пытается реплицировать только те страницы, которые активно используются на разных NUMA-нодах во время выполнения приложения. В основе реализации данной политики лежит существующий механизм автоматической NUMA-балансировки Linux. Фаза сканирования виртуального адресного пространства NUMA-балансировщиком остаётся практически без изменений. Единственное отличие заключается в том, что реплицируемые виртуальные области памяти не исключаются из рассмотрения независимо от типа отображения~--- включая \textit{разделяемые file-backed} отображения.

Дальнейшие изменения происходят уже на этапе обработки \textit{NUMA hinting fault}. В стандартной реализации Linux при обнаружении удалённого доступа к памяти страница может быть мигрирована на NUMA-узел, с которого выполняется обращение. В предлагаемом механизме вместо миграции выполняется репликация страницы. К этому моменту необходимые таблицы трансляций уже гарантированно существуют в реплицированном виде, что позволяет корректно разместить локальные копии данных на соответствующих NUMA-нодах.

\subsubsection{Дерепликация таблиц и данных}
\label{sec:Chapter3_data_table_derepl}

Помимо механизмов репликации, реализация также предусматривает выполнение обратного процесса~--- дерепликации таблиц трансляций и данных. Необходимость в этом возникает при изменении свойств виртуальных областей памяти или смене политики репликации процесса на менее агрессивную.

Дерепликация таблиц трансляций требуется только при ослаблении политики репликации таблиц. Например, при переходе с \texttt{All} на \texttt{Minimal} часть ранее реплицированных структур становится избыточной и должна быть возвращена в обычное состояние. Аналогичная ситуация возникает и при полном отключении репликации.

Дерепликация данных выполняется в двух случаях:
\begin{enumerate}
    \item При отключении репликации данных для процесса.
    \item При изменении прав доступа к области памяти посредством системного вызова \texttt{mprotect}, если \textit{read-only} отображение становится \textit{writable}. Поскольку механизм репликации данных рассчитан исключительно на \textit{read-only} страницы, перед предоставлением возможности записи данные должны быть предварительно возвращены в обычное нереплицированное состояние.
\end{enumerate}

Важно, что дерепликация данных всегда должна выполняться раньше дерепликации таблиц.

Также стоит отметить, что при дерепликации данных ни одна из существующих реплик не переиспользуется в качестве основной страницы. Вместо этого аллоцируется новая физическая страница, в которую копируется содержимое одной из реплик, после чего остальные экземпляры освобождаются. Причины такого решения во многом аналогичны тем, что рассматривались ранее при обсуждении отказа от переиспользования исходной страницы во время репликации данных:
\begin{enumerate}
    \item Во-первых, реплицированные страницы являются намеренно упрощёнными с точки зрения интеграции с подсистемами \textit{reverse mapping} и \textit{swap}. Аллокация новой страницы автоматически возвращает память в стандартное состояние с полноценным участием во всех механизмах управления памятью ядра.

    \item Во-вторых, не существует однозначного критерия выбора одной из реплик в качестве оригинальной страницы. При создании нового экземпляра страница выделяется на NUMA-ноде, наиболее подходящей с точки зрения действующей политики распределения памяти процесса.

    \item Наконец, наиболее важная причина связана со всё тем же совместным использованием памяти между процессами. Как и в случае с обычными \textit{анонимными private} страницами, реплицированные страницы могут \textit{разделяться} между несколькими процессами после выполнения \texttt{clone}. Переиспользование одной из реплик привело бы к ситуации, когда одна и та же физическая страница одновременно рассматривается как реплицированная в одном процессе и как обычная~--- в другом, что нарушило бы согласованность механизма репликации (см главу~\ref{sec:Chapter4_phys_layer}).
\end{enumerate}

С таблицами трансляций описанных выше проблем не возникает, поэтому при их дерепликации переиспользуется один из экземпляров.

\subsection{Поддержка больших страниц}
\label{sec:Chapter3_huge_pages}

Одной из задач работы являлась поддержка \textit{больших страниц}, в частности механизмов \textit{Transparent Huge Pages}\cite{THP} и \textit{HugeTLB}\cite{hugetlb}. На практике оказалось, что их реализация содержит большое количество важных архитектурных особенностей и тонкостей, которые необходимо учитывать при сочетании с механизмом репликации таблиц трансляций и данных. Далее рассматриваются ключевые из них.

\subsubsection{\textit{compound}-страницы, \textit{reverse mapping} и \texttt{deferred\_split\_folio}}

В Linux \textit{большие страницы} реализуются как последовательность подряд идущих физических страниц базового размера, объединённых в единую логическую сущность~--- \textit{compound}-страницу (или \textit{folio} в более современной терминологии ядра). Первая базовая страница выступает в роли \textit{head}-страницы и хранит общие свойства всей \textit{compound}-страницы (например, большинство флагов состояния), тогда как остальные являются \textit{tail}-страницами и могут содержать локальную для своей базовой страницы информацию.

На рис.~\ref{fig:compound_pages} показаны два варианта отображения \textit{compound}-страницы в виртуальное адресное пространство процесса. В первом случае отображение осуществляется через PMD-запись, покрывающую всю большую страницу целиком. Во втором случае та же \textit{compound}-страница отображается набором обычных PTE-записей.

\begin{figure}[H]
    \centering
    \includegraphics[width=1\textwidth]{compound_page.pdf}
    \caption{Отображение \textit{compound}-страницы через PMD- и PTE-записи}
    \label{fig:compound_pages}
\end{figure}

Ввиду возможности отображать \textit{compound}-страницы через PTE-записи могут возникать ситуации, когда \textit{compound}-страница отображена лишь частично. Например, после выполнения системного вызова \texttt{munmap} часть диапазона виртуальных адресов может перестать ссылаться на соответствующие физические страницы, тогда как оставшиеся участки памяти продолжат использовать ту же \textit{compound}-страницу; освободившийся диапазон виртуальных адресов впоследствии может быть занят другой областью памяти, использующей уже иные физические страницы. Аналогичная ситуация может возникнуть и в результате операции \texttt{MADV\_DONTNEED}, когда часть физических базовых страниц отбрасывается, а новые страницы при последующих обращениях будут выделены уже вне исходной \textit{compound}-страницы. И т.д. (см рис.~\ref{fig:compound_partial_mapping}).

\begin{figure}[H]
    \centering
    \includegraphics[width=1\textwidth]{compound_partial_mapping.pdf}
    \caption{Частичное использование \textit{compound}-страниц}
    \label{fig:compound_partial_mapping}
\end{figure}

Проблема заключается в том, что \textit{compound}-страница остаётся единой логической сущностью. Даже если часть входящих в неё базовых страниц фактически больше не используется, они всё равно продолжают удерживаться в составе \textit{compound}-страницы и не могут быть освобождены независимо друг от друга.

Эту проблему решает механизм \textit{reverse mapping}. Он позволяет отслеживать фактическое использование базовых страниц в составе большой страницы и выявлять ситуации, когда некоторые из них перестают быть используемыми. Такие \textit{compound}-страницы помещаются в очередь \texttt{deferred\_split\_folio}, предназначенную для последующего отложенного разделения \textit{compound}-страницы на набор независимых базовых страниц. После разделения неиспользуемые базовые страницы уже могут быть корректно освобождены системой, особенно в условиях нехватки памяти.

Что касается механизма репликации данных, то было принято решение не поддерживать \texttt{deferred\_split\_folio} для реплицированных \textit{compound}-страниц по двум причинам:
\begin{enumerate}
    \item \textit{reverse mapping} не поддерживается для реплицированных данных согласно дизайну. Изменение этого решения на столь поздней стадии разработки потребовало бы существенной переработки реализации.
    \item \textit{reverse mapping}~--- лишь одна из многочисленных тонкостей реализации \texttt{deferred\_split\_folio}. Помимо этого потребовалось ещё много с чем разбираться, что привело бы к дальнейшему усложнению и без того непростого дизайна.
\end{enumerate}

\subsubsection{\texttt{khugepaged}}

Одним из ключевых механизмов поддержки \textit{Transparent Huge Pages} в Linux является фоновый поток ядра \texttt{khugepaged}. Его задача заключается в периодическом сканировании адресных пространств процессов и поиске диапазонов виртуальной памяти, которые могут быть преобразованы (свёрнуты) из набора обычных страниц в одну большую.  При этом исходное состояние памяти может быть достаточно неоднородным: некоторые PTE-записи могут отсутствовать вовсе, ссылаться на \texttt{zero\_page} или находиться в \textit{swap}-состоянии. Существующие же физические страницы могут быть выделены независимо друг от друга и располагаться на разных NUMA-нодах. В случае успешной свёртки такой набор страниц заменяется новой \textit{compound}-страницей, размещённой на одной NUMA-ноде и отображаемой через PMD-запись (см рис.~\ref{fig:khugepaged_collapse}).

Важно уточнить, что речь идёт прежде всего об \textit{анонимной} памяти: в рассматриваемой версии ядра поддержка свёртки для \textit{file-backed} памяти была в экспериментальном состоянии.

\begin{figure}[H]
    \centering
    \includegraphics[width=1\textwidth]{khugepaged.pdf}
    \caption{Свёртка набора базовых страниц в PMD-mapped \textit{compound}-страницу}
    \label{fig:khugepaged_collapse}
\end{figure}

При интеграции с механизмом репликации для виртуальных областей памяти с реплицируемыми данными свёртка была запрещена. Теоретически их поддержку можно было бы реализовать, предварительно выполняя дерепликацию данных перед свёрткой, однако практическая целесообразность такого подхода вызывает сомнения. Есть опасение, что механизм репликации данных, с одной стороны, будет создавать реплики страниц, а \texttt{khugepaged}, с другой,~--- практически сразу устранять их посредством дерепликации. Это могло бы привести к круговороту лишней работы и ухудшению производительности. Поэтому было принято решение не совмещать оба механизма.

Тем не менее, реплицированные PMD-mapped \textit{compound}-страницы всё же могут возникнуть. Например, большая страница может быть первоначально выделена обработчиком ошибок страниц, а затем позднее реплицирована посредством \texttt{phys\_duplicate}. Для этого функция \texttt{phys\_duplicate} была модифицирована и получила поддержку репликации PMD-mapped \textit{compound}-страниц.

Поддержка реплицированных таблиц трансляций при этом всё же является необходимой. В противном случае включение механизма репликации фактически парализовало бы работу \texttt{khugepaged}, а задача по интеграции репликации в \textit{Transparent Huge Pages} была бы не выполнена. С точки зрения репликации таблиц трансляций возможны три случая:
\begin{enumerate}
    \item PMD-слой не реплицирован, PTE-слой также не реплицирован~--- стандартная реализация Linux;
    \item PMD-слой реплицирован, PTE-слой~--- нет;
    \item\label{itm:khugepaged} PMD- и PTE-слои реплицированы, данные при этом не реплицированы.
\end{enumerate}

Рассмотрим случай~\ref{itm:khugepaged}, поскольку он является наиболее общим и архитектурно сложным:
\begin{enumerate}
    \item\label{itm:khugepaged_pre} Перед началом свёртки выполняются подкачка \textit{swap}-страниц и аллокация будущей \textit{compound}-страницы, чтобы избежать возможных проблем с этим на более поздних этапах операции.

    \item Все PMD-записи, соответствующие репликам PTE-таблиц, временно удаляются из таблиц трансляций, но их исходные значения сохраняются (см рис.~\ref{fig:collapse_1}). Выполняется очистка соответствующего диапазона адресов из TLB-кеша всех процессоров.

    Это необходимо, во-первых, для удаления из TLB-кеша записей, относящихся к страницам, участвующим в свёртке, а во-вторых, для предотвращения обращений пользовательских потоков к подменяемым физическим страницам как через PTE-таблицы, так и через уже закешированные TLB-записи.

    Таким образом, если попытка обращения всё же произойдёт, возникнет ошибка страниц, обработка которой будет невозможна до завершения свёртки, поскольку операция выполняется при эксклюзивном захвате соответствующей виртуальной области памяти.

    При этом сохранение исходных PMD-записей принципиально важно: каждая из них ссылается на собственную PTE-таблицу, и в случае неудачи все эти таблицы потребуется восстановить. Если же свёртка завершится успешно, лишние PTE-таблицы будут освобождены (см пункт~\ref{itm:deposit}).

    \begin{figure}[H]
        \centering
        \includegraphics[width=1\textwidth]{collapse_1.pdf}
        \caption{Временное удаление PMD-записей во время свёртки}
        \label{fig:collapse_1}
    \end{figure}

    \item Далее \texttt{khugepaged} сканирует изолированные PTE-записи диапазона памяти и выполняет различные проверки корректности свёртки. Если выясняется, что свёртка более невозможна, выполняется откат с восстановлением сохранённых ранее PMD-записей.

    \item В случае успеха содержимое исходных страниц копируется в \textit{compound}-страницу, аллоцированную на этапе~\ref{itm:khugepaged_pre}. Затем старые PTE-записи всех реплик очищаются, а соответствующие базовые страницы освобождаются (см рис.~\ref{fig:collapse_2}).

    \begin{figure}[H]
        \centering
        \includegraphics[width=1\textwidth]{collapse_2.pdf}
        \caption{Копирование данных в новую \textit{compound}-страницу}
        \label{fig:collapse_2}
    \end{figure}

    \item\label{itm:deposit} На последнем этапе новая PMD-mapped \textit{compound}-страница устанавливается в записи всех реплик таблиц трансляций.

    Кроме того, одна из более ненужных PTE-таблиц сохраняется для возможного использования в будущем. Дело в том, что PMD-mapped отображение большой страницы впоследствии может потребоваться разбить обратно на набор PTE-mapped отображений (подробнее об этом будет сказано в главе~\ref{sec:Chapter3_split_huge_pmd}). Чтобы в момент такого разбиения гарантированно не возникло проблем с аллокацией новой PTE-таблицы, существует список отложенных (\textit{deposited}) PTE-таблиц, куда и помещается одна из освободившихся реплик. А т.к. сама большая страница данных не реплицирована, достаточно сохранить лишь одну PTE-таблицу, тогда как остальные реплики освобождаются (см рис.~\ref{fig:collapse_3}).

    \begin{figure}[H]
        \centering
        \includegraphics[width=1\textwidth]{collapse_3.pdf}
        \caption{Установка \textit{compound}-страницы в PMD-записи}
        \label{fig:collapse_3}
    \end{figure}
\end{enumerate}

\subsubsection{\texttt{split\_huge\_pmd}}
\label{sec:Chapter3_split_huge_pmd}

Несмотря на то, что большие страницы в первую очередь предназначены для отображения в виртуальное адресное пространство через PMD-записи, в ряде ситуаций ядру требуется работать лишь с отдельными базовыми страницами в составе \textit{compound}-страницы. Например, это необходимо при выполнении системных вызовов \texttt{munmap} или \texttt{madvise} с флагом \texttt{MADV\_DONTNEED}, когда требуется удалить отображение или освободить только часть диапазона памяти. Аналогично, \texttt{mprotect} может изменять права доступа лишь для части большой страницы, а \texttt{mremap}~--- перемещать отдельные диапазоны виртуальной памяти. Такие операции возможны потому, что базовый размер страницы продолжает оставаться минимальной единицей отображения памяти как на уровне физической памяти, так и на уровне виртуального адресного пространства процесса.

Для поддержки подобных операций в Linux используется функция \texttt{split\_huge\_pmd}. Её задача заключается в преобразовании PMD-mapped отображения \textit{compound}-страницы в набор обычных PTE-mapped отображений (см рис.~\ref{fig:compound_pages}). При этом сама \textit{compound}-страница продолжает существовать как единая логическая сущность: изменяется лишь способ её отображения в таблицах трансляций.

При интеграции с механизмом репликации возможны три случая:
\begin{enumerate}
    \item PMD-слой не реплицирован, \textit{compound}-страница тоже~--- стандартная реализация Linux;
    \item PMD-слой реплицирован, \textit{compound}-страница~--- нет;
    \item\label{itm:split_huge_pmd} PMD-слой и \textit{compound}-страница реплицированы.
\end{enumerate}

Рассмотрим случай~\ref{itm:split_huge_pmd} как архитектурно наиболее сложный:
\begin{enumerate}
    \item Все PMD-записи, соответствующие репликам \textit{compound}-страницы, временно инвалидируются, после чего выполняется очистка соответствующего диапазона адресов из TLB-кеша процессоров.

    Это необходимо для того, чтобы пользовательские потоки не могли обращаться к PMD-mapped отображению в момент перестройки. В отличие от \texttt{khugepaged}, где взаимное исключение обеспечивалось эксклюзивным захватом виртуальной области памяти, здесь оно достигается посредством захвата \textit{page table lock}, общего для всех реплик PMD-таблицы (см рис.~\ref{fig:split_huge_pmd_1}).

    \begin{figure}[H]
        \centering
        \includegraphics[width=1\textwidth]{split_huge_pmd_1.pdf}
        \caption{Инвалидация PMD-записей перед разбиением}
        \label{fig:split_huge_pmd_1}
    \end{figure}

    \item Далее из списка \textit{deposited} PTE-таблиц извлекается по одной таблице для каждой реплики \textit{compound}-страницы. Необходимое количество таблиц гарантированно присутствует в списке, поскольку они были туда заранее сохранены функцией \texttt{phys\_duplicate} во время репликации большой страницы.

    Все записи новых PTE-таблиц инициализируются таким образом, чтобы каждая из них указывала на соответствующую базовую страницу внутри своей реплики \textit{compound}-страницы (см рис.~\ref{fig:split_huge_pmd_2}).

    \begin{figure}[H]
        \centering
        \includegraphics[width=1\textwidth]{split_huge_pmd_2.pdf}
        \caption{Подготовка PTE-таблиц}
        \label{fig:split_huge_pmd_2}
    \end{figure}

    \item На последнем этапе подготовленные PTE-таблицы встраиваются в иерархии таблиц трансляций: каждая таблица подключается к соответствующей реплике PMD-таблицы вместо ранее инвалидированной PMD-записи (см рис.~\ref{fig:split_huge_pmd_3}).

    \begin{figure}[H]
        \centering
        \includegraphics[width=1\textwidth]{split_huge_pmd_3.pdf}
        \caption{Установка PTE-таблиц в реплики PMD-таблиц}
        \label{fig:split_huge_pmd_3}
    \end{figure}
\end{enumerate}

\subsubsection{Contiguous PTE ranges}
\label{sec:Chapter3_contpte}

Следует напомнить, что механизм репликации разрабатывался не для исходного Linux 6.6, а для дистрибутива openEuler 25.03. Он хоть и основан на Linux 6.6, часть функциональностей была портирована в него из более поздних версий Linux, вплоть до Linux 6.11. В частности, в openEuler присутствует механизм \textit{multi-size Transparent Huge Pages} (mTHP)\cite{mTHP}, появившийся в Linux 6.9.

В отличие от классических \textit{Transparent Huge Pages}, где большая страница имеет фиксированный размер и отображается одной PMD-записью, mTHP позволяет использовать большие страницы меньших размеров, являющихся степенями двойки относительно базового размера страницы (см рис.~\ref{fig:mTHP}). Такие страницы по-прежнему представляют собой \textit{compound}-страницы, однако отображаются только через набор PTE-записей. При этом виртуальный и физический адреса должны быть выровнены в соответствии с размером отображаемой страницы.

\begin{figure}[H]
    \centering
    \includegraphics[width=1\textwidth]{mTHP.pdf}
    \caption{multi-size THP}
    \label{fig:mTHP}
\end{figure}

Вместе с mTHP в ядре появилась функциональность \textit{contiguous PTE ranges}\cite{contpte} для архитектуры arm64~--- именно в ней кроется главное преимущество PTE-mapped \textit{compound}-страниц по сравнению с обычной последовательностью подряд идущих PTE-mapped базовых страниц. Её идея заключается в том, что в arm64 все PTE-записи разбиты на \textit{contiguous} диапазоны по 16 записей в каждом (см рис.~\ref{fig:contpte}). Если все записи в составе диапазона ссылаются на последовательные физические страницы и имеют одинаковые атрибуты, то диапазон можно \textit{свернуть}, установив \textit{contig} бит во всех входящих в него PTE-записях.

\begin{figure}[H]
    \centering
    
\includegraphics[width=1\textwidth]{contpte.pdf}
    \caption{Contiguous PTE ranges}
    \label{fig:contpte}
\end{figure}

Это позволит TLB-кешу использовать одну запись для всего \textit{contiguous} диапазона вместо шестнадцати, что приведёт к более эффективному его потреблению и в конечном счёте уменьшит количество TLB-промахов. А возможность формировать такие диапазоны как раз и обеспечивается механизмом \textit{compound}-страниц, гарантирующим непрерывность физических страниц в памяти.

Кроме того, работа с \textit{compound}-страницами ускоряет выполнение некоторых операций управления памятью, таких как \texttt{clone}, \texttt{munmap} и \texttt{madvise} с флагом \texttt{MADV\_DONTNEED}, поскольку ядро получает возможность обрабатывать PTE-записи \textit{compound}-страниц пачками вместо их поэлементной обработки.

Однако поддержка \textit{contiguous PTE ranges} существенно усложняет реализацию функций, модифицирующих PTE-записи. Теперь перед изменением отдельной записи ядру может потребоваться предварительно \textit{разбить} \textit{contiguous} диапазон, сбросив \textit{contig} бит у PTE-записей в нём. А после модификации множества подряд идущих записей~--- попытаться \textit{свернуть} их в \textit{contiguous range}.

В итоге именно поддержка \textit{contiguous PTE ranges} стала наиболее сложной частью интеграции больших страниц с механизмом репликации. Несмотря на то, что значительная часть инфраструктуры поддержки больших страниц уже была реализована, завершить поддержку \textit{contiguous PTE ranges} в поставленные сроки не удалось.
